% ── §1  Introduction ──────────────────────────────────────────────────────────
\section{Introduction}
\label{sec:intro}

Distributed streaming pipelines increasingly support applications with
hard real-time constraints: real-time financial fraud detection, autonomous
vehicle perception stacks, and sub-second interactive analytics.
In these systems, \emph{tail latency}---the \pnn{99} or \pnn{999}
end-to-end processing delay---rather than average throughput determines
whether a deadline is met~\cite{dean2013tail}.
A single slow request that crosses a 500\ms SLA boundary is a failure
regardless of what the mean latency is.

Modern stream processing engines, Apache Flink~\cite{carbone2015flink}
and Apache Kafka~\cite{kreps2011kafka} chief among them, have rich knob
spaces: checkpoint intervals govern fault-tolerance overhead;
buffer timeouts trade batching efficiency for added latency;
operator parallelism determines how evenly load is distributed.
Yet these knobs are set \emph{statically} at deployment time.
When load spikes, a TaskManager fails, or a noisy neighbour contends for
CPU, the configuration that was optimal at startup becomes harmful.
Production operators respond with manual re-tuning or coarse autoscalers
that react too slowly and overshoot~\cite{welsh2001seda}.

We identify three structural gaps in the state of the art:
\begin{enumerate}
  \item \textbf{No cross-layer view.} Tail spikes emerge from interactions
    across the ingest layer (Kafka consumer lag), the execution layer
    (Flink buffer backpressure), and the state layer (checkpoint stalls).
    Existing autoscalers address at most one layer.
  \item \textbf{No closed-loop learning.} Rule-based threshold systems
    react but do not generalise. They require hand-crafted thresholds that
    vary by workload and cluster configuration.
  \item \textbf{No explainability.} Even if a controller works, operators
    need to understand its decisions to trust it in production.
\end{enumerate}

We present \xtal (\textbf{eX}plainable \textbf{T}ail-latency
\textbf{A}daptive contro\textbf{L}ler), a system that closes all three
gaps simultaneously.
\xtal runs a 10-second control loop that: (i)~collects a six-dimensional
telemetry snapshot from Prometheus; (ii)~selects a joint Flink+Kafka
configuration from a discrete action space using an XGBoost contextual
bandit trained online; (iii)~applies the configuration through real APIs;
and (iv)~produces per-decision SHAP feature attributions that show which
telemetry signal drove the control decision.
A hard safety layer enforces correctness invariants (e.g., checkpoint
interval $\geq 500$\ms) independently of the learned policy.

\paragraph{Contributions.}
\begin{itemize}
  \item \xtal, the first closed-loop, cross-layer, explainable
    tail-latency controller for a real Kafka\,+\,Flink stack~(\S\ref{sec:design}).
  \item A baseline characterisation study of where tail spikes originate
    under steady, burst, and failure conditions~(\S\ref{sec:motivation}).
  \item An online XGBoost contextual bandit formulation with a provably
    safe action space and SHAP explainability, requiring no offline
    training~(\S\ref{sec:design}).
  \item An evaluation on NEXMark benchmarks showing up to 38\% \pnn{99}
    reduction and 7$\times$ lower SLA violation rate versus the
    next-best baseline~(\S\ref{sec:eval}).
\end{itemize}
